% theme: information_communication_network
\MajorQuestion{情報通信ネットワーク}
\SetRowSpaceQanda

\Instruction{次の問いに答えなさい。}
\QQ{情報を盗む、機器を破壊するなど、悪意のある不正プログラムを何というか。}{マルウェア}
\QQ{家庭や学校など、限られた範囲のネットワークを何というか。}{LAN}
\QQ{ネットワーク上で情報機器を識別するために割り当てられる数字の組み合わせを\NamedBlank{ip}という。}{IPアドレス}
\QQ{ネットワークを通じて、コンピュータやソフトウェアを利用できるサービスシステムを何というか。}{クラウドコンピューティング}
\QQ{ネットワークどうしを中継し、データの交通整理を行う機器を何というか。}{ルータ}

\Instruction{\strut}
\QQ{機器の故障などに備えて、データの予備を保存しておくことを何というか。}{バックアップ}
\QQ{世界規模のネットワークを何というか。}{インターネット}
\QQ{ネットワークを安全に利用するための技術を\NamedBlank{security}という。}{情報セキュリティ}
\QQ{インターネットで、データを効率よく送信するために分割した小さなまとまりを何というか。}{パケット}
\QQ{人間には分かりにくい\RefBlank{ip}に対応させる、文字列の名前を\NamedBlank{domain}という。}{ドメイン名}

\Instruction{\strut}
\QQ{許可された通信のみを内部に通すしくみを何というか。}{ファイアウォール}
\QQ{インターネットにつながるための標準的な通信方式を何というか。}{TCP/IP}
\QQ{\RefBlank{ip}と\RefBlank{domain}を対応させるサーバを何というか。}{DNSサーバ}
\QQ{ネットワーク上で、情報やサービスを提供するコンピュータを何というか。}{サーバ}
\QQ{ブラウザと相手のコンピュータ間の通信の暗号化技術として使われるものを何というか。}{SSL}

\Instruction{\strut}
\QQ{情報機器が通信するときに従う必要がある約束を何というか。}{ネットワークプロトコル}
\QQ{情報の漏えいを防ぐために、データを第三者に読みにくい形に変えることを何というか。}{暗号化}
\QQ{情報機器どうしを有線や無線でつなぎ、情報のやりとりができるようにしたものを何というか。}{情報通信ネットワーク}
\QQ{仮想空間における\RefBlank{security}を何というか。}{サイバーセキュリティ}
\QQ{コンピュータウイルスなどへの対策に利用するソフトウェアを何というか。}{セキュリティ対策ソフトウェア}
